\subsection*{一步一步建立直观}

\paragraph{1. 相对变化给出自然的局部单位}

若存量为 $a>0$，很小的绝对增加 $h$ 对应相对增加约 $h/a$。自然对数把这种相对增量变成输出增量，所以局部斜率是 $1/a$。右图在 $a=1$ 附近的切线斜率为一，正好给出一套没有额外比例常数的刻度。



\paragraph{2. 不是每个函数都测量相对变化}

切换到平方函数，比较同样的观察位置与步长。它的斜率是 $2a$，随存量增大而增大，与对数恰好相反。这一步用于突出自然对数的特殊尺度；右图切换函数后的读数属于新函数，不能把它继续解释成对数导数。



\paragraph{3. 从导数回到积累的解释}

回到对数曲线，试着解释从一增长到 $x$ 的对数值：它等于沿途累计 $1/t$ 的面积。因为每段按当前存量归一化，同样的比例增长产生相同的对数增量。将起点与终点保持为正，面积解释便与对数的乘法规律互相支持。


